\documentclass[twocolumn]{aastex631}
\begin{document}
\title{The reionization ionizing-photon-budget shortfall for star-forming galaxies is not robust to systematics at z\~{}6 (beta systematic corner)}
\author{NebulaMind Lab (autonomous pipeline)}
\affiliation{NebulaMind Lab, \url{https://lab.nebulamind.net}}
\begin{abstract}
Reionization ionizing-photon-budget reconciliation at z\~{}6: star-forming galaxies require f\_esc=0.048 (+0.048/-0.025) to close the budget (Madau-Dickinson SFRD, log xi\_ion=25.5+/-0.15, clumping C=2-5, JWST-SFRD tail), versus indirect-proxy-inferred f\_esc=0.050 (+0.076/-0.030) from LzLCS O32/beta calibrations. Median delta(required-inferred)=-0.002 dex-frac (16-84\%: -0.075 to +0.050); 48\% of the systematic MC shows a shortfall. Conclusion: the budget CLOSES within the systematic, and the sign holds under both O32 and beta calibrations. The result is bounded by the xi\_ion x clumping x proxy-calibration systematic, not by statistics. Generated autonomously from public data (jwst) via the NebulaMind Lab runner: a bounded, reproducible, descriptive study.
\end{abstract}
\section{Introduction}
Recent studies have highlighted a potential crisis in reconciling the ionizing photon budget during reionization [Muñoz2024]. This has led to increased demands on ionizing sources, with some suggesting that absorption-dominated reionization may require additional contributions from other sources [Davies2021]. To address this issue, we revisit the ionizing photon budget using a literature-anchored approach. Our work builds upon previous efforts to calibrate excursion set reionization models and assess the galaxy ionizing photon budget at high redshifts [Park2022, Duncan2015].
\section{Data and method}
In our analysis, we utilize the cosmic star formation rate density (SFRD) from Madau \& Dickinson's analytic fitting function [Madau2017]. We adopt published values for xi\_ion and O32/beta f\_esc proxy calibrations from LzLCS studies [Chisholm+22, Flury+22; Simmonds+24]. By combining these elements, we calculate the ionizing photon budget using a systematic approach that accounts for uncertainties in clumping factor (C) and escape fraction (f\_esc).
\section{Result}
Our reconciliation of the reionization ionizing-photon-budget at z\~{}6 reveals that star-forming galaxies require an escape fraction f\_esc = 0.048 (+0.048/-0.025) to close the budget, assuming a Madau-Dickinson SFRD, log xi\_ion=25.5+/-0.15, and clumping factor C between 2-5. This value is consistent with indirect-proxy-inferred f\_esc = 0.050 (+0.076/-0.030) from LzLCS O32/beta calibrations. The median delta between required and inferred values is -0.002 dex-frac, ranging from -0.075 to +0.050 (16-84\% confidence interval), with 48\% of systematic Monte Carlo simulations showing a shortfall.
\begin{figure}\centering\includegraphics[width=\columnwidth]{result.png}\caption{The reionization ionizing-photon-budget shortfall for star-forming galaxies is not robust to systematics at z\~{}6 (beta systematic corner)}\end{figure}
\section{Caveats}
Despite these findings, it is essential to acknowledge the limitations of our approach. Our analysis relies on automated, single-selection, and uncalibrated measurements, which may introduce biases and uncertainties. The accuracy of our results depends heavily on the adopted values for xi\_ion, clumping factor, and proxy calibrations from previous studies. Furthermore, our method does not account for potential contributions from other ionizing sources or additional systematic errors in the input data. Therefore, while our study provides a valuable reconciliation of the reionization photon budget, further research is needed to refine these estimates and address the underlying uncertainties.
\section*{References}
\footnotesize
[Muoz2024] Muñoz et al. (2024). Reionization after JWST: a photon budget crisis?. bibcode 2024MNRAS.535L..37M \par [Davies2021] Davies et al. (2021). The Predicament of Absorption-dominated Reionization: Increased Demands on Ionizing Sources. bibcode 2021ApJ...918L..35D \par [Park2022] Park et al. (2022). Calibrating excursion set reionization models to approximately conserve ionizing photons. bibcode 2022MNRAS.517..192P \par [Duncan2015] Duncan et al. (2015). Powering reionization: assessing the galaxy ionizing photon budget at z < 10. bibcode 2015MNRAS.451.2030D \par [Madau2017] Madau (2017). Cosmic Reionization after Planck and before JWST: An Analytic Approach. bibcode 2017ApJ...851...50M

\end{document}
